\documentclass[12pt]{article}
\usepackage[utf8]{inputenc}
\usepackage[T1]{fontenc}
\usepackage{hyperref}
\title{Zhuang Language Resources for K-12 Curriculum Development}
\author{Compiled by Claude Sonnet 5.5 (claude-sonnet-5-5) for Andr\'as Kornai}
\date{}
\begin{document}
\maketitle

\section{Introduction}
The code zha is the ISO 639-3 macrolanguage code for Zhuang. Standard Zhuang (Vahcuengh) is the official standardized form, based on Northern Zhuang, and has used a Latin-based script since 1982. It has six tones. The searches returned reference and overview material. They found no graded school textbook with a title and publisher.

\section{Primary Educational Resources}
\begin{itemize}
\item A mobile app for learning the romanized Zhuang writing system was listed. The search results do not say what it covers beyond that. \url{https://apps.apple.com/app/id1369853980}
\end{itemize}
No school textbooks or primers with verifiable details were found.

\section{Linguistic and Reference Materials}
\begin{itemize}
\item Wikipedia, ``Standard Zhuang''. It cites reference works such as the 2006 \textit{Northern Zhuang-Chinese-Thai-English Dictionary} by Somsonge Burusphat and Qin Xiaohang (Bangkok). \url{https://en.wikipedia.org/wiki/Standard_Zhuang}
\item Omniglot, Zhuang writing systems. \url{https://omniglot.com/writing/zhuang.htm}
\item Zhuang alphabet article. \url{https://demo.azizisearch.com/lite/wikipedia/page/Zhuang_alphabet}
\item Unicode IRG document on Zhuang characters (kZhuang). \url{https://www.unicode.org/irg/docs/n2677-kZhuang.pdf}
\item Journal article on Zhuang, Thai Journal Online (full details not verified). \url{https://so03.tci-thaijo.org/index.php/JLC/article/download/21647/18702/46569}
\item Wikipedia, ISO 639 entry on the Zhuang languages. \url{https://wikipedia.com/wiki/ISO_639:ccx}
\end{itemize}

\section{Religious and Community Resources}
None were found in the search results.

\section{Digital and Audio Resources}
Only the app listed above was found. No audio or video resources were found.

\section{Curriculum Development Context}
Zhuang has an official standard orthography and a proficiency test (VSSG) at three levels, as reported in the search summary. Further searches in Chinese for Zhuang-medium school textbooks, such as those published in Guangxi, would be the next step.

\end{document}